\documentclass[11pt,a4paper]{article}
\usepackage[utf8]{inputenc}
\usepackage[T1]{fontenc}
\usepackage[french]{babel}
\usepackage{geometry}
\usepackage{graphicx}
\usepackage{xcolor}
\usepackage{listings}
\usepackage{hyperref}
\usepackage{fancyhdr}
\usepackage{titlesec}
\usepackage{tcolorbox}
\usepackage{booktabs}
\usepackage{array}
\usepackage{tabularx}
\usepackage{enumitem}
\usepackage{fontawesome5}

\geometry{margin=2.5cm, top=3cm, bottom=2.5cm}

% ============================================================================
% COULEURS PKF TELCO
% ============================================================================
\definecolor{pkfblue}{RGB}{31,56,100}
\definecolor{pkforange}{RGB}{236,130,54}
\definecolor{pkfgray}{RGB}{80,80,80}
\definecolor{codebg}{RGB}{248,248,248}
\definecolor{codeframe}{RGB}{220,220,220}
\definecolor{successgreen}{RGB}{40,167,69}
\definecolor{warningyellow}{RGB}{255,193,7}
\definecolor{dangered}{RGB}{220,53,69}

% ============================================================================
% STYLES DES TITRES
% ============================================================================
\titleformat{\section}{\Large\bfseries\color{pkfblue}}{\thesection}{1em}{}[\color{pkforange}\titlerule]
\titleformat{\subsection}{\large\bfseries\color{pkfblue}}{\thesubsection}{1em}{}
\titleformat{\subsubsection}{\normalsize\bfseries\color{pkfgray}}{\thesubsubsection}{1em}{}

\titlespacing*{\section}{0pt}{20pt}{10pt}
\titlespacing*{\subsection}{0pt}{15pt}{8pt}

% ============================================================================
% STYLE DU CODE
% ============================================================================
\lstdefinestyle{bashstyle}{
    language=bash,
    backgroundcolor=\color{codebg},
    basicstyle=\ttfamily\small,
    breaklines=true,
    frame=single,
    rulecolor=\color{codeframe},
    numbers=left,
    numberstyle=\tiny\color{gray},
    keywordstyle=\color{pkfblue}\bfseries,
    commentstyle=\color{successgreen},
    stringstyle=\color{pkforange},
    showstringspaces=false,
    tabsize=2,
    xleftmargin=1.5em,
    framexleftmargin=1.5em,
    aboveskip=1em,
    belowskip=1em
}

\lstdefinestyle{yamlstyle}{
    backgroundcolor=\color{codebg},
    basicstyle=\ttfamily\small,
    breaklines=true,
    frame=single,
    rulecolor=\color{codeframe},
    numbers=left,
    numberstyle=\tiny\color{gray},
    keywordstyle=\color{pkfblue},
    stringstyle=\color{pkforange},
    showstringspaces=false,
    tabsize=2,
    xleftmargin=1.5em,
    framexleftmargin=1.5em,
    aboveskip=1em,
    belowskip=1em
}

\lstdefinestyle{confstyle}{
    backgroundcolor=\color{codebg},
    basicstyle=\ttfamily\small,
    breaklines=true,
    frame=single,
    rulecolor=\color{codeframe},
    numbers=left,
    numberstyle=\tiny\color{gray},
    showstringspaces=false,
    tabsize=2,
    xleftmargin=1.5em,
    framexleftmargin=1.5em,
    aboveskip=1em,
    belowskip=1em
}

\lstset{style=bashstyle}

% ============================================================================
% EN-TÊTE ET PIED DE PAGE
% ============================================================================
\pagestyle{fancy}
\fancyhf{}
\fancyhead[L]{\textcolor{pkfblue}{\textbf{PKF Telco}}}
\fancyhead[C]{\textcolor{pkfgray}{\small Guide CMAS/PWS v1.0}}
\fancyhead[R]{\textcolor{pkfblue}{\thepage}}
\fancyfoot[C]{\textcolor{pkfgray}{\small Cameroun — Octobre 2026}}
\renewcommand{\headrulewidth}{2pt}
\renewcommand{\headrule}{\hbox to\headwidth{\color{pkforange}\leaders\hrule height \headrulewidth\hfill}}
\renewcommand{\footrulewidth}{0.5pt}
\renewcommand{\footrule}{\hbox to\headwidth{\color{codeframe}\leaders\hrule height \footrulewidth\hfill}}

% ============================================================================
% BOÎTES PERSONNALISÉES
% ============================================================================
\tcbuselibrary{skins,breakable}

% Boîte pour les commandes
\newtcolorbox{cmdbox}[1][]{
    enhanced,
    colback=codebg,
    colframe=pkfblue,
    coltitle=white,
    fonttitle=\bfseries\small,
    title=#1,
    boxrule=1pt,
    arc=2pt,
    left=5pt,
    right=5pt,
    top=5pt,
    bottom=5pt,
    breakable
}

% Boîte pour les avertissements
\newtcolorbox{warnbox}{
    enhanced,
    colback=warningyellow!15,
    colframe=pkforange,
    coltitle=white,
    fonttitle=\bfseries\small,
    title={\faExclamationTriangle\ \ Important},
    boxrule=1pt,
    arc=2pt,
    left=5pt,
    right=5pt,
    breakable
}

% Boîte pour les succès
\newtcolorbox{successbox}{
    enhanced,
    colback=successgreen!10,
    colframe=successgreen,
    coltitle=white,
    fonttitle=\bfseries\small,
    title={\faCheckCircle\ \ Vérification},
    boxrule=1pt,
    arc=2pt,
    left=5pt,
    right=5pt,
    breakable
}

% Boîte pour les erreurs
\newtcolorbox{errorbox}{
    enhanced,
    colback=dangered!10,
    colframe=dangered,
    coltitle=white,
    fonttitle=\bfseries\small,
    title={\faTimesCircle\ \ Attention},
    boxrule=1pt,
    arc=2pt,
    left=5pt,
    right=5pt,
    breakable
}

% Boîte pour les fichiers
\newtcolorbox{filebox}[1][]{
    enhanced,
    colback=codebg,
    colframe=pkfgray,
    coltitle=white,
    fonttitle=\bfseries\ttfamily\small,
    title={\faFile\ \ #1},
    boxrule=1pt,
    arc=2pt,
    left=0pt,
    right=0pt,
    top=0pt,
    bottom=0pt,
    breakable
}

% ============================================================================
% COMMANDES PERSONNALISÉES
% ============================================================================
\newcommand{\ip}[1]{\texttt{\color{pkfblue}#1}}
\newcommand{\port}[1]{\texttt{\color{pkforange}#1}}
\newcommand{\cmd}[1]{\texttt{\color{pkfgray}#1}}
\newcommand{\file}[1]{\texttt{\color{successgreen}#1}}

% ============================================================================
% DÉBUT DU DOCUMENT
% ============================================================================
\begin{document}

% ============================================================================
% PAGE DE TITRE
% ============================================================================
\begin{titlepage}
    \centering
    \vspace*{1cm}
    
    % Logo/Titre
    {\fontsize{50}{60}\selectfont\textcolor{pkfblue}{\textbf{PKF}}}\\[0.2cm]
    {\fontsize{30}{36}\selectfont\textcolor{pkforange}{\textbf{Telco}}}\\[1.5cm]
    
    % Ligne décorative
    {\color{pkforange}\rule{0.6\textwidth}{3pt}}\\[1.5cm]
    
    % Titre principal
    {\Huge\bfseries\textcolor{pkfblue}{Guide d'Installation}}\\[0.5cm]
    {\Huge\bfseries\textcolor{pkfblue}{Système CMAS/PWS}}\\[1cm]
    
    {\Large\textcolor{pkfgray}{Cell Broadcast Emergency Alerts}}\\[0.3cm]
    {\Large\textcolor{pkfgray}{pour Réseau LTE Privé}}\\[2cm]
    
    % Technologies
    \begin{tcolorbox}[
        enhanced,
        colback=codebg,
        colframe=pkfblue,
        width=0.7\textwidth,
        arc=5pt,
        boxrule=1pt
    ]
        \centering
        \textbf{Stack Technique}\\[0.3cm]
        \textcolor{pkfblue}{Open5GS} + \textcolor{pkforange}{srsRAN 4G} + \textcolor{successgreen}{osmo-cbc}
    \end{tcolorbox}
    
    \vfill
    
    % Informations
    \begin{tabular}{rl}
        \textcolor{pkfgray}{\textbf{Version}} & 1.0 \\[0.2cm]
        \textcolor{pkfgray}{\textbf{Date}} & Octobre 2026 \\[0.2cm]
        \textcolor{pkfgray}{\textbf{Auteur}} & Arnaud DJOUM \\[0.2cm]
        \textcolor{pkfgray}{\textbf{Organisation}} & F2G / PKF Telco \\
    \end{tabular}
    
    \vspace{1cm}
    
    % Ligne décorative bas
    {\color{pkforange}\rule{0.8\textwidth}{2pt}}
    
\end{titlepage}

% ============================================================================
% TABLE DES MATIÈRES
% ============================================================================
\tableofcontents
\newpage

% ============================================================================
% SECTION 1: INTRODUCTION
% ============================================================================
\section{Introduction}

Ce document décrit l'installation complète d'un système d'alertes d'urgence \textbf{CMAS} (Commercial Mobile Alert System) / \textbf{PWS} (Public Warning System) sur un réseau LTE privé PKF Telco.

\subsection{Architecture du système}

\begin{center}
\begin{tabular}{|>{\columncolor{pkfblue!10}}l|l|l|}
\hline
\rowcolor{pkfblue}
\textcolor{white}{\textbf{Composant}} & \textcolor{white}{\textbf{Rôle}} & \textcolor{white}{\textbf{Interface}} \\
\hline
osmo-cbc & Cell Broadcast Centre & REST API + SBc-AP \\
\hline
Open5GS MME & Mobility Management Entity & SBc-AP + S1AP \\
\hline
srsRAN eNB & Station de base & S1AP + Air (SIB12) \\
\hline
UE & Téléphone mobile & Interface radio \\
\hline
\end{tabular}
\end{center}

\subsection{Flux des alertes}

\begin{center}
\begin{tcolorbox}[
    enhanced,
    colback=codebg,
    colframe=pkfblue,
    width=0.9\textwidth,
    arc=3pt
]
\centering
\texttt{API REST} $\rightarrow$ \textbf{osmo-cbc} $\xrightarrow{\text{SBc-AP}}$ \textbf{MME} $\xrightarrow{\text{S1AP}}$ \textbf{eNB} $\xrightarrow{\text{SIB12}}$ \textbf{UE}
\end{tcolorbox}
\end{center}

\subsection{Prérequis matériels et logiciels}

\begin{itemize}[leftmargin=2cm]
    \item[\faServer] Ubuntu 22.04 LTS (ou supérieur)
    \item[\faDocker] Docker et Docker Compose
    \item[\faBroadcastTower] bladeRF xA4 (ou SDR compatible srsRAN)
    \item[\faSimCard] Carte SIM programmable (sysmoISIM-SJA2)
    \item[\faCode] Git, CMake, GCC, build-essential
\end{itemize}

\subsection{Adresses IP du déploiement}

\begin{center}
\begin{tabular}{|l|l|l|}
\hline
\rowcolor{pkfblue}
\textcolor{white}{\textbf{Composant}} & \textcolor{white}{\textbf{Adresse IP}} & \textcolor{white}{\textbf{Port}} \\
\hline
MongoDB & \ip{172.22.0.2} & \port{27017} \\
\hline
HSS & \ip{172.22.0.3} & \port{3868} (Diameter) \\
\hline
SGW-C & \ip{172.22.0.5} & \port{2123} (GTP-C) \\
\hline
SGW-U & \ip{172.22.0.6} & \port{2152} (GTP-U) \\
\hline
SMF & \ip{172.22.0.7} & \port{2123} \\
\hline
UPF & \ip{172.22.0.8} & \port{2152} \\
\hline
MME (GTP) & \ip{172.22.0.1} & \port{2123} \\
\hline
MME (S1AP) & \ip{127.0.0.2} & \port{36412} \\
\hline
MME (SBc-AP) & \ip{127.0.0.2} & \port{29168} \\
\hline
osmo-cbc (API) & \ip{127.0.0.1} & \port{12345} \\
\hline
osmo-cbc (SBc-AP) & \ip{127.0.0.10} & \port{29169} \\
\hline
eNB (S1-C) & \ip{127.0.1.1} & - \\
\hline
\end{tabular}
\end{center}

\newpage

% ============================================================================
% SECTION 2: INSTALLATION OPEN5GS
% ============================================================================
\section{Installation Open5GS avec SBc-AP}

\begin{warnbox}
La version standard d'Open5GS ne supporte pas l'interface SBc-AP. Il faut compiler depuis la Pull Request \#4864 qui ajoute ce support.
\end{warnbox}

\subsection{Clonage du code source}

\begin{cmdbox}[Étape 2.1 — Clonage du repository]
\begin{lstlisting}[style=bashstyle]
# Créer le répertoire de travail
mkdir -p /home/f2g/Desktop
cd /home/f2g/Desktop

# Cloner Open5GS
git clone https://github.com/open5gs/open5gs.git open5gs-fork
cd open5gs-fork

# Récupérer la branche SBc-AP (PR #4864)
git fetch origin pull/4864/head:sbcap-branch
git checkout sbcap-branch
\end{lstlisting}
\end{cmdbox}

\subsection{Installation des dépendances}

\begin{cmdbox}[Étape 2.2 — Dépendances système]
\begin{lstlisting}[style=bashstyle]
sudo apt update
sudo apt install -y \
    meson ninja-build flex bison \
    libsctp-dev libgnutls28-dev libgcrypt-dev \
    libssl-dev libidn11-dev libmongoc-dev \
    libbson-dev libyaml-dev libnghttp2-dev \
    libmicrohttpd-dev libcurl4-gnutls-dev \
    libjansson-dev libtalloc-dev libtins-dev
\end{lstlisting}
\end{cmdbox}

\subsection{Compilation et installation}

\begin{cmdbox}[Étape 2.3 — Compilation]
\begin{lstlisting}[style=bashstyle]
cd /home/f2g/Desktop/open5gs-fork

# Configuration avec Meson
meson build --prefix=/usr/local

# Compilation
ninja -C build

# Installation (nécessite sudo)
sudo ninja -C build install
\end{lstlisting}
\end{cmdbox}

\subsection{Vérification de l'installation}

\begin{successbox}
\begin{lstlisting}[style=bashstyle]
# Vérifier que le binaire MME supporte SBc-AP
strings /usr/local/bin/open5gs-mmed | grep -i sbcap

# Résultat attendu:
# sbcap_server
# SBCAP accepted
# SBc-AP
\end{lstlisting}
\end{successbox}

\newpage

% ============================================================================
% SECTION 3: INSTALLATION OSMO-CBC
% ============================================================================
\section{Installation osmo-cbc}

osmo-cbc est le Cell Broadcast Centre développé par Osmocom. Il gère les messages d'alerte et communique avec le MME via SBc-AP.

\subsection{Installation des dépendances Osmocom}

\begin{cmdbox}[Étape 3.1 — Dépendances Osmocom]
\begin{lstlisting}[style=bashstyle]
sudo apt install -y \
    libosmocore-dev libosmo-netif-dev \
    libosmo-sigtran-dev libulfius-dev \
    libjansson-dev autoconf automake libtool
\end{lstlisting}
\end{cmdbox}

\subsection{Compilation depuis les sources}

\begin{cmdbox}[Étape 3.2 — Compilation osmo-cbc]
\begin{lstlisting}[style=bashstyle]
cd /home/f2g
git clone https://gitea.osmocom.org/cellular-infrastructure/osmo-cbc
cd osmo-cbc

# Génération des scripts de configuration
autoreconf -fi

# Configuration
./configure --prefix=/usr/local

# Compilation
make -j$(nproc)

# Installation
sudo make install
\end{lstlisting}
\end{cmdbox}

\subsection{Vérification}

\begin{successbox}
\begin{lstlisting}[style=bashstyle]
# Vérifier l'installation
which osmo-cbc
# /usr/local/bin/osmo-cbc

osmo-cbc --version
\end{lstlisting}
\end{successbox}

\begin{warnbox}
Si vous avez un binaire patché pour la gestion de l'expiration des alertes, utilisez:
\texttt{/usr/bin/osmo-cbc-patched}
\end{warnbox}

\newpage

% ============================================================================
% SECTION 4: INSTALLATION SRSRAN
% ============================================================================
\section{Installation srsRAN avec support PWS/CMAS}

\begin{errorbox}
La version standard de srsRAN 4G \textbf{ne génère pas} le SIB12 pour les alertes CMAS. Le code affiche uniquement \texttt{"would generate SIB12"} sans réellement le diffuser. Il faut utiliser la branche \texttt{pws-cmas-support} et s'assurer que le code génère réellement le SIB12.
\end{errorbox}

\subsection{Clonage et sélection de la branche}

\begin{cmdbox}[Étape 4.1 — Clonage srsRAN]
\begin{lstlisting}[style=bashstyle]
mkdir -p /home/f2g/build
cd /home/f2g/build

git clone https://github.com/srsran/srsRAN_4G.git
cd srsRAN_4G

# Checkout de la branche PWS
git checkout pws-cmas-support
\end{lstlisting}
\end{cmdbox}

\subsection{Vérification du code source}

\begin{warnbox}
Avant de compiler, vérifiez que le fichier \file{srsenb/src/stack/rrc/rrc.cc} contient bien:
\begin{center}
\texttt{logger.info("PWS: CMAS warning - generating SIB12 for broadcast");}
\end{center}
Et \textbf{PAS}:
\begin{center}
\texttt{logger.info("PWS: CMAS warning - would generate SIB12");}
\end{center}
\end{warnbox}

\begin{cmdbox}[Étape 4.2 — Vérification du code]
\begin{lstlisting}[style=bashstyle]
grep -n "generating SIB12\|would generate" \
    /home/f2g/build/srsRAN_4G/srsenb/src/stack/rrc/rrc.cc
\end{lstlisting}
\end{cmdbox}

\subsection{Installation des dépendances}

\begin{cmdbox}[Étape 4.3 — Dépendances srsRAN]
\begin{lstlisting}[style=bashstyle]
sudo apt install -y \
    cmake libfftw3-dev libmbedtls-dev \
    libboost-program-options-dev libconfig++-dev \
    libsctp-dev libuhd-dev libbladerf-dev \
    libsoapysdr-dev
\end{lstlisting}
\end{cmdbox}

\subsection{Compilation}

\begin{cmdbox}[Étape 4.4 — Compilation srsRAN]
\begin{lstlisting}[style=bashstyle]
cd /home/f2g/build/srsRAN_4G
mkdir -p build && cd build

cmake .. -DCMAKE_BUILD_TYPE=Release
make -j$(nproc)

sudo make install
sudo ldconfig
\end{lstlisting}
\end{cmdbox}

\subsection{Vérification du binaire}

\begin{successbox}
\begin{lstlisting}[style=bashstyle]
# Vérifier que le binaire génère bien SIB12
strings $(which srsenb) | grep "generating SIB12"

# Résultat attendu:
# PWS: CMAS warning - generating SIB12 for broadcast
\end{lstlisting}
\end{successbox}

\newpage

% ============================================================================
% SECTION 5: CONFIGURATION DOCKER
% ============================================================================
\section{Configuration du Core Network Docker}

Le cœur de réseau (HSS, SGW, SMF, UPF) fonctionne dans des conteneurs Docker.

\subsection{Structure des conteneurs}

\begin{center}
\begin{tabular}{|l|l|l|l|}
\hline
\rowcolor{pkfblue}
\textcolor{white}{\textbf{Conteneur}} & \textcolor{white}{\textbf{IP}} & \textcolor{white}{\textbf{Rôle}} & \textcolor{white}{\textbf{Interface}} \\
\hline
mongo & \ip{172.22.0.2} & Base de données & MongoDB \\
\hline
hss & \ip{172.22.0.3} & Home Subscriber Server & Diameter S6a \\
\hline
sgwc & \ip{172.22.0.5} & SGW Control Plane & GTP-C \\
\hline
sgwu & \ip{172.22.0.6} & SGW User Plane & GTP-U \\
\hline
smf & \ip{172.22.0.7} & Session Management & GTP-C \\
\hline
upf & \ip{172.22.0.8} & User Plane Function & GTP-U \\
\hline
\end{tabular}
\end{center}

\subsection{Démarrage des conteneurs}

\begin{cmdbox}[Étape 5.1 — Démarrage Docker]
\begin{lstlisting}[style=bashstyle]
cd /home/f2g/docker_open5gs

# Démarrer les conteneurs nécessaires
docker-compose up -d mongo hss sgwc sgwu smf upf

# Vérifier le statut
docker-compose ps
\end{lstlisting}
\end{cmdbox}

\subsection{Vérification de la connectivité}

\begin{successbox}
\begin{lstlisting}[style=bashstyle]
# Vérifier que le HSS est accessible
ping -c 2 172.22.0.3

# Vérifier que le SGW-C est accessible
ping -c 2 172.22.0.5
\end{lstlisting}
\end{successbox}

\newpage

% ============================================================================
% SECTION 6: CONFIGURATION MME
% ============================================================================
\section{Configuration du MME avec SBc-AP}

\subsection{Création de la structure de répertoires}

\begin{cmdbox}[Étape 6.1 — Création des répertoires]
\begin{lstlisting}[style=bashstyle]
mkdir -p /tmp/f2g_cmas/{mme,cbc,enb1,certs}
\end{lstlisting}
\end{cmdbox}

\subsection{Fichier de configuration MME}

\begin{filebox}[/tmp/f2g\_cmas/mme/mme-sbcap.yaml]
\begin{lstlisting}[style=yamlstyle]
logger:
  file:
    path: /tmp/f2g_cmas/mme.log
  level: info
  domain: mme,sbcap,s1ap,sbc

mme:
  freeDiameter: /tmp/mme-fork-fd.conf
  s1ap:
    server:
      - address: 127.0.0.2
        port: 36412
      - address: 192.168.1.100
        port: 36412
  sbcap:
    server:
      - address: 127.0.0.2
        port: 29168
  gtpc:
    server:
      - address: 172.22.0.1
    client:
      sgwc:
        - address: 172.22.0.5
      smf:
        - address: 172.22.0.7
  gummei:
    - plmn_id:
        mcc: 001
        mnc: 01
      mme_gid: 2
      mme_code: 1
  tai:
    - plmn_id:
        mcc: 001
        mnc: 01
      tac: 1
    - plmn_id:
        mcc: 001
        mnc: 01
      tac: 2
  security:
    integrity_order: [EIA2, EIA1, EIA0]
    ciphering_order: [EEA0, EEA1, EEA2]
  network_name:
    full: PKF Telco
  mme_name: mme-pkf-cmas
\end{lstlisting}
\end{filebox}

\newpage

\subsection{Configuration freeDiameter}

\begin{filebox}[/tmp/mme-fork-fd.conf]
\begin{lstlisting}[style=confstyle]
Identity = "mme.epc.mnc001.mcc001.3gppnetwork.org";
Realm = "epc.mnc001.mcc001.3gppnetwork.org";
Port = 3868;
SecPort = 5868;
No_SCTP;

ConnectPeer = "hss.epc.mnc001.mcc001.3gppnetwork.org" { 
    ConnectTo = "172.22.0.3"; 
    Port = 3868; 
    No_TLS; 
};

LoadExtension = "/usr/local/lib/x86_64-linux-gnu/freeDiameter/dict_s6a.fdx";
LoadExtension = "/usr/local/lib/x86_64-linux-gnu/freeDiameter/dict_nas_mipv6.fdx";
\end{lstlisting}
\end{filebox}

\begin{warnbox}
L'adresse IP \ip{172.22.0.1} pour \texttt{gtpc.server} est l'IP du bridge Docker. C'est \textbf{critique} pour que les conteneurs Docker puissent atteindre le MME.
\end{warnbox}

\newpage

% ============================================================================
% SECTION 7: CONFIGURATION OSMO-CBC
% ============================================================================
\section{Configuration osmo-cbc}

\begin{filebox}[/tmp/f2g\_cmas/cbc/osmo-cbc.cfg]
\begin{lstlisting}[style=confstyle]
log stderr
 logging color 1
 logging print category 1
 logging timestamp 1
 logging level all info
 logging level sbcap debug
 logging level rest debug

cbc
 unknown-peers reject
 
 ecbe
  local-ip 127.0.0.1
  local-port 12345
 
 sbcap
  local-ip 127.0.0.10
  local-port 29169
 
 peer sbcap pkf-mme
  mode client
  remote-ip 127.0.0.2
  remote-port 29168
\end{lstlisting}
\end{filebox}

\begin{center}
\begin{tabular}{|l|l|}
\hline
\rowcolor{pkfblue}
\textcolor{white}{\textbf{Paramètre}} & \textcolor{white}{\textbf{Description}} \\
\hline
\texttt{ecbe.local-port} & Port de l'API REST (12345) \\
\hline
\texttt{sbcap.local-ip} & IP locale pour SBc-AP \\
\hline
\texttt{peer.remote-ip} & IP du MME \\
\hline
\texttt{peer.remote-port} & Port SBc-AP du MME (29168) \\
\hline
\end{tabular}
\end{center}

\newpage

% ============================================================================
% SECTION 8: CONFIGURATION ENB
% ============================================================================
\section{Configuration eNB (srsRAN)}

\subsection{Configuration principale}

\begin{filebox}[/tmp/f2g\_cmas/enb1/enb.conf]
\begin{lstlisting}[style=confstyle]
[enb]
enb_id = 0x19B01
mcc = 001
mnc = 01
mme_addr = 127.0.0.2
gtp_bind_addr = 172.22.0.1
s1c_bind_addr = 127.0.1.1
s1c_bind_port = 0
n_prb = 50

[enb_files]
sib_config = /tmp/f2g_cmas/sib_pws.conf
rr_config = /home/f2g/.config/srsran/rr.conf
rb_config = /home/f2g/.config/srsran/rb.conf

[rf]
device_name = bladeRF
device_args = auto
time_adv_nsamples = auto
tx_gain = 60
rx_gain = 40

[log]
filename = /tmp/f2g_cmas/enb1.log
all_level = info
phy_level = warning
\end{lstlisting}
\end{filebox}

\newpage

\subsection{Configuration SIB avec SIB12}

\begin{errorbox}
Le \textbf{SIB12 DOIT} être déclaré dans \texttt{si\_mapping\_info = [ 12 ]} pour que l'eNB puisse diffuser les alertes CMAS. Sans cette configuration, les alertes sont ignorées.
\end{errorbox}

\begin{filebox}[/tmp/f2g\_cmas/sib\_pws.conf]
\begin{lstlisting}[style=confstyle]
sib1 =
{
    intra_freq_reselection = "Allowed";
    q_rx_lev_min = -65;
    cell_barred = "NotBarred"
    si_window_length = 20;
    sched_info =
    (
        {
            si_periodicity = 16;
            si_mapping_info = [ 3 ];
        },
        {
            si_periodicity = 64;
            si_mapping_info = [ 12 ];  // CRITIQUE: SIB12 pour CMAS
        }
    );
    system_info_value_tag = 11;
};

sib2 = 
{
    rr_config_common_sib =
    {
        rach_cnfg = { 
            num_ra_preambles = 52; 
            preamble_init_rx_target_pwr = -104;
            pwr_ramping_step = 6;
            preamble_trans_max = 10;
            ra_resp_win_size = 10;
            mac_con_res_timer = 64;
            max_harq_msg3_tx = 4;
        };
        bcch_cnfg = { modification_period_coeff = 16; };
        pcch_cnfg = { default_paging_cycle = 32; nB = "1"; };
        prach_cnfg = { 
            root_sequence_index = 128;
            prach_cnfg_info = {
                high_speed_flag = false;
                prach_config_index = 3;
                prach_freq_offset = 4;
                zero_correlation_zone_config = 5;
            };
        };
        pdsch_cnfg = { p_b = 1; rs_power = 0; };
        pusch_cnfg = { 
            n_sb = 1;
            hopping_mode = "inter-subframe";
            pusch_hopping_offset = 2;
            enable_64_qam = false;
            ul_rs = {
                cyclic_shift = 0;
                group_assignment_pusch = 0;
                group_hopping_enabled = false;
                sequence_hopping_enabled = false;
            };
        };
        pucch_cnfg = { 
            delta_pucch_shift = 1;
            n_rb_cqi = 1;
            n_cs_an = 0;
            n1_pucch_an = 12;
        };
        ul_pwr_ctrl = { 
            p0_nominal_pusch = -85;
            alpha = 0.7;
            p0_nominal_pucch = -107;
            delta_flist_pucch = {
                format_1 = 0; format_1b = 3;
                format_2 = 1; format_2a = 2; format_2b = 2;
            };
            delta_preamble_msg3 = 6;
        };
        ul_cp_length = "len1";
    };
    ue_timers_and_constants = {
        t300 = 2000; t301 = 100; t310 = 200;
        n310 = 1; t311 = 10000; n311 = 1;
    };
    freqInfo = {
        ul_carrier_freq_present = true;
        ul_bw_present = true;
        additional_spectrum_emission = 1;
    };
    time_alignment_timer = "INFINITY";
};

sib3 =
{
    q_hyst = 2;
    threshServingLow = 2;
};
\end{lstlisting}
\end{filebox}

\newpage

% ============================================================================
% SECTION 9: SCRIPT D'ALERTES
% ============================================================================
\section{Script d'envoi d'alertes}

\begin{filebox}[/tmp/alert.sh]
\begin{lstlisting}[style=bashstyle]
#!/bin/bash
# Script d'envoi d'alertes CMAS pour PKF Telco
CBC_API="http://127.0.0.1:12345/api/v1"

case "$1" in
  send)
    MSG="$2"
    MSG_ID="${3:-4370}"
    DURATION="${4:-60}"
    SCOPE="${5:-tac:1}"
    
    CODE=$((RANDOM % 100 + 1))
    NAME="alert${MSG_ID}c${CODE}"
    EXPIRE=$((DURATION + 5))
    
    echo "Sending: $NAME (msg_id=$MSG_ID, code=$CODE)"
    echo "Duration: ${DURATION}s - Expiration: ${EXPIRE}s"
    echo "Scope: $SCOPE"
    
    MSG_HEX=$(echo -n "$MSG" | xxd -p | tr -d '\n')
    
    if [[ "$SCOPE" == tac:* ]]; then
      TAC="${SCOPE#tac:}"
      SCOPE_JSON="\"scope_tai\": {\"tai\": [{\"plmn\": \"00101\", \"tac\": $TAC}]}"
    else
      SCOPE_JSON="\"scope_plmn\": {\"plmn\": \"00101\"}"
    fi
    
    RESULT=$(curl -s -X POST "$CBC_API/message" \
      -H "Content-Type: application/json" \
      -d "{
        \"cbe_name\": \"$NAME\",
        \"category\": \"cmas\",
        \"repetition_period\": 5,
        \"num_of_bcast\": $((DURATION / 5)),
        \"dcs\": 15,
        \"msg_id\": $MSG_ID,
        \"payload_encoded\": \"$MSG_HEX\",
        $SCOPE_JSON
      }" 2>&1)
    
    if echo "$RESULT" | grep -qi "201\|Created\|success"; then
      echo "Result: Created"
    else
      echo "Result: $RESULT"
    fi
    ;;
    
  stop)
    MSG_ID="${2:-4370}"
    echo "Stopping msg_id=$MSG_ID..."
    curl -s -X DELETE "$CBC_API/message/$MSG_ID"
    echo "Done"
    ;;
    
  list)
    echo "Active messages:"
    curl -s "$CBC_API/message" | python3 -m json.tool 2>/dev/null || echo "None"
    ;;
    
  *)
    echo "Usage: $0 {send|stop|list}"
    echo "  send \"message\" [msg_id] [duration_sec] [tac:N]"
    echo "  stop [msg_id]"
    echo "  list"
    ;;
esac
\end{lstlisting}
\end{filebox}

\begin{cmdbox}[Rendre le script exécutable]
\begin{lstlisting}[style=bashstyle]
chmod +x /tmp/alert.sh
\end{lstlisting}
\end{cmdbox}

\newpage

% ============================================================================
% SECTION 10: DÉMARRAGE
% ============================================================================
\section{Démarrage du système}

\subsection{Ordre de démarrage}

\begin{center}
\begin{tabular}{|c|l|l|}
\hline
\rowcolor{pkfblue}
\textcolor{white}{\textbf{Étape}} & \textcolor{white}{\textbf{Composant}} & \textcolor{white}{\textbf{Dépendances}} \\
\hline
1 & Docker (HSS, SGW, SMF, UPF) & - \\
\hline
2 & MME Open5GS & Docker \\
\hline
3 & osmo-cbc & MME \\
\hline
4 & eNB srsRAN & MME, bladeRF \\
\hline
\end{tabular}
\end{center}

\subsection{Commandes de démarrage}

\begin{cmdbox}[Étape 10.1 — Démarrage Docker]
\begin{lstlisting}[style=bashstyle]
cd /home/f2g/docker_open5gs
docker-compose up -d mongo hss sgwc sgwu smf upf
\end{lstlisting}
\end{cmdbox}

\begin{cmdbox}[Étape 10.2 — Démarrage MME]
\begin{lstlisting}[style=bashstyle]
/usr/local/bin/open5gs-mmed -c /tmp/f2g_cmas/mme/mme-sbcap.yaml &
\end{lstlisting}
\end{cmdbox}

\begin{cmdbox}[Étape 10.3 — Démarrage osmo-cbc]
\begin{lstlisting}[style=bashstyle]
cd /tmp/f2g_cmas/cbc
/usr/bin/osmo-cbc-patched -c osmo-cbc.cfg &
# ou: /usr/local/bin/osmo-cbc -c osmo-cbc.cfg &
\end{lstlisting}
\end{cmdbox}

\begin{cmdbox}[Étape 10.4 — Démarrage eNB]
\begin{lstlisting}[style=bashstyle]
cd /tmp/f2g_cmas/enb1
sudo srsenb enb.conf &
\end{lstlisting}
\end{cmdbox}

\subsection{Vérification des connexions}

\begin{successbox}
\begin{lstlisting}[style=bashstyle]
# Vérifier les logs MME
grep -E "CONNECTED|SBC/CBC|Number of eNBs" /tmp/f2g_cmas/mme.log

# Résultat attendu:
# CONNECTED TO 'hss.epc.mnc001.mcc001.3gppnetwork.org'
# SBC/CBC accepted[127.0.0.10]
# [Added] Number of eNBs is now 1
\end{lstlisting}
\end{successbox}

\newpage

% ============================================================================
% SECTION 11: UTILISATION
% ============================================================================
\section{Envoi d'alertes}

\subsection{Exemples d'utilisation}

\begin{cmdbox}[Envoyer une alerte sur TAC 1]
\begin{lstlisting}[style=bashstyle]
/tmp/alert.sh send "URGENCE INONDATION DOUALA" 4370 60 tac:1
\end{lstlisting}
\end{cmdbox}

\begin{cmdbox}[Envoyer une alerte sur TAC 2]
\begin{lstlisting}[style=bashstyle]
/tmp/alert.sh send "ALERTE SEISME YAOUNDE" 4370 120 tac:2
\end{lstlisting}
\end{cmdbox}

\begin{cmdbox}[Lister les alertes actives]
\begin{lstlisting}[style=bashstyle]
/tmp/alert.sh list
\end{lstlisting}
\end{cmdbox}

\begin{cmdbox}[Stopper une alerte]
\begin{lstlisting}[style=bashstyle]
/tmp/alert.sh stop 4370
\end{lstlisting}
\end{cmdbox}

\subsection{Vérification de la diffusion}

\begin{successbox}
\begin{lstlisting}[style=bashstyle]
# Vérifier les logs eNB pour la diffusion SIB12
grep -iE "PWS|SIB12|CMAS" /tmp/f2g_cmas/enb1.log | tail -10

# Résultat attendu:
# PWS WriteReplaceWarning: msg_id=0x1112, serial=0x4120, CMAS
# PWS: CMAS warning - generating SIB12 for broadcast
# PWS: SIB12 cc=0 SI#2 segment 1/1 written (96 bytes)
# PWS paging: cmas indication scheduled in 32 paging occasions
# PWS: CMAS Alert Broadcast - msg_id=0x1112 serial=0x4120
\end{lstlisting}
\end{successbox}

\newpage

% ============================================================================
% SECTION 12: DÉPANNAGE
% ============================================================================
\section{Dépannage}

\subsection{Problèmes courants et solutions}

\begin{center}
\begin{tabularx}{\textwidth}{|l|l|X|}
\hline
\rowcolor{pkfblue}
\textcolor{white}{\textbf{Symptôme}} & \textcolor{white}{\textbf{Cause}} & \textcolor{white}{\textbf{Solution}} \\
\hline
"Could not allocate" & Alerte précédente active & \cmd{/tmp/alert.sh stop 4370} ou redémarrer osmo-cbc \\
\hline
"would generate SIB12" & Binaire non à jour & Recompiler srsRAN avec la branche pws-cmas-support \\
\hline
UE ne s'attache pas & GTP IP incorrecte & Utiliser \ip{172.22.0.1} pour gtpc.server \\
\hline
HSS non connecté & Mauvaise IP/realm & Vérifier \file{/tmp/mme-fork-fd.conf} \\
\hline
eNB ne démarre pas & bladeRF non détecté & Débrancher/rebrancher USB \\
\hline
SIB12 non diffusé & SIB12 non configuré & Ajouter \texttt{si\_mapping\_info = [ 12 ]} dans sib.conf \\
\hline
\end{tabularx}
\end{center}

\subsection{Commandes de diagnostic}

\begin{cmdbox}[Diagnostic système]
\begin{lstlisting}[style=bashstyle]
# État du bladeRF
bladeRF-cli -e "version"

# Logs MME en temps réel
tail -f /tmp/f2g_cmas/mme.log | grep -E "Attach|IMSI|Bearer"

# Logs eNB en temps réel
tail -f /tmp/f2g_cmas/enb1.log | grep -iE "PWS|SIB12"

# Processus actifs
pgrep -a "open5gs-mmed\|osmo-cbc\|srsenb"

# Connexions réseau
netstat -an | grep -E "29168|36412|12345"
\end{lstlisting}
\end{cmdbox}

\newpage

% ============================================================================
% SECTION 13: RÉFÉRENCES
% ============================================================================
\section{Références}

\subsection{Spécifications 3GPP}

\begin{itemize}[leftmargin=2cm]
    \item[\faBook] \textbf{3GPP TS 23.041} — Technical realization of Cell Broadcast Service (CBS)
    \item[\faBook] \textbf{3GPP TS 29.168} — Cell Broadcast Centre interfaces (SBc-AP)
    \item[\faBook] \textbf{3GPP TS 36.331} — E-UTRA RRC Protocol (SIB10/11/12)
    \item[\faBook] \textbf{3GPP TS 36.413} — S1 Application Protocol (S1AP)
\end{itemize}

\subsection{Documentation des projets}

\begin{itemize}[leftmargin=2cm]
    \item[\faGithub] \textbf{Open5GS} — \url{https://open5gs.org/open5gs/docs/}
    \item[\faGithub] \textbf{Open5GS PR \#4864} — SBc-AP Support
    \item[\faGithub] \textbf{srsRAN 4G} — \url{https://docs.srsran.com/projects/4g}
    \item[\faGithub] \textbf{osmo-cbc} — \url{https://osmocom.org/projects/osmo-cbc}
\end{itemize}

\subsection{Message Identifiers (3GPP TS 23.041)}

\begin{center}
\begin{tabular}{|l|l|l|}
\hline
\rowcolor{pkfblue}
\textcolor{white}{\textbf{Message ID}} & \textcolor{white}{\textbf{Type}} & \textcolor{white}{\textbf{Description}} \\
\hline
4370 (0x1112) & Presidential Alert & Alerte présidentielle \\
\hline
4371-4372 & Extreme Alert & Menace extrême \\
\hline
4373-4378 & Severe Alert & Menace sévère \\
\hline
4379-4392 & Amber Alert & Alerte enlèvement \\
\hline
4393-4394 & Public Safety & Sécurité publique \\
\hline
4395 & State/Local Test & Test local \\
\hline
\end{tabular}
\end{center}

\vfill

% ============================================================================
% PIED DE PAGE FINAL
% ============================================================================
\begin{center}
{\color{pkforange}\rule{0.8\textwidth}{2pt}}\\[0.5cm]
{\Large\textbf{\textcolor{pkfblue}{PKF Telco}}}\\[0.3cm]
\textcolor{pkfgray}{Cameroun — Système CMAS/PWS}\\[0.3cm]
\textcolor{pkfgray}{Document généré le \today}\\[0.3cm]
\textcolor{pkfgray}{\small Version 1.0 — Arnaud DJOUM / F2G}
\end{center}

\end{document}
